% !TEX root = ../../../main.tex
% 所属：第一部分：单变量分析学 / 第3章：实数的八大基本定理
% 节：3.1 拓扑空间与开闭集
% 正文按原文件顺序保存；由 ../chapter3.tex 调用。
\section{拓扑空间与开闭集}
\begin{definition}\label{defoftosp}
	令$X$是一个集合.$X$的拓扑$\mathcal{T}$是由$X$的一些子集所构成的族，且$\mathcal{T}$满足
	\begin{enumerate}
		\item $X\in \mathcal{T}$且$\varnothing\in\mathcal{T}$
		\item $\mathcal{T}$ 中集合的任意并仍属于 $\mathcal{T}$
		\item $\mathcal{T}$中集合的有限交仍属于$\mathcal{T}$
	\end{enumerate}
	称偶对$\left(\mathrm{X}, \mathcal{T}\right)$ 是一个拓扑空间.
\end{definition}
\begin{definition}
	$\left(\mathrm{X}, \mathcal{T}\right)$ 是一个拓扑空间,那么$\mathcal{T}$中的元素在$\mathrm{X}$中称做开的，换言之$\mathcal{T}$中的元素称为$\mathrm{X}$中的开集.
	
	闭集定义为开集关于$\X$的补集.
	
	在$\left(\mathrm{X}, \mathcal{T}\right)$ 中存在即是开集又是闭集的集，称为闭开集.
\end{definition}
\begin{proposition}
	对拓扑空间$\left(\mathrm{X}, \mathcal{T}\right)$
	
	（1）$\mathrm{X}$ ，$\varnothing$ 是团开的
	
	（2）开集的任意并仍是开集
	
	（3）开集的有限交仍是开集
	
	（4）闭集的任意交仍是闭集
	
	（5）闭集的有限并仍是闭集
\end{proposition}
\begin{proof}
	$(1),(2),(3)$由开集的定义是显然的，对$(4),(5)$使用De Morgan定理,设$S_{\alpha}\in\mathcal{T}$,那么$S_{\alpha}^C$是闭集，则
	$$\bigcap_\alpha S_\alpha^C=\left(\bigcup_\alpha^n S_\alpha\right)^C \text {, 而 }\bigcup_\alpha^n S_\alpha  \in \mathcal{T}$$
	$$	\bigcup_{i=1}^n S_i^C=\left(\bigcap_{i=1}^n S_i\right)^C \text {, 而 } \bigcap_{i=1}^n S_i \in \mathcal{T}$$
\end{proof}
\begin{example}
	\begin{enumerate}
		\item		$\mathrm{X}=  \{a, b, c, d, e, f\}$,令
		$$\mathcal{T}_1=\{\mathrm{X}, \phi,\{a\},\{c, d\},\{a, c, d\},\{b, c, d, e, f\}]$$
		显然地，$\mathcal{T}_1$是$\mathrm{X}$上的一个拓扑：其中
		\begin{itemize}
			\item $\{a\}$是闭开的
			\item$\{b,c\}$既不是闭的，也不是开的是开的，而不是闭的
			\item$\{c,d\}$
			\item$\{a, b, e, f\}$ 是闭的，而不是开的
		\end{itemize}
		\item 离散拓扑$\T$,对集合$\mathrm{X}$,$\T$定义为$\mathrm{X}$所有子集构成的族.显然的，$\mathrm{X}$的任意子集都是闭开的.
		\item 平凡拓扑$\T_0$,定义为
		\[\T_0:=\{\mathrm{X},\varnothing\}\]
	\end{enumerate}
\end{example}
\begin{theorem}
	$f:X\rightarrow Y$,是一个从$X$到$Y$的$1-1$映射，且是满射，令$\left(\mathrm{X}, \mathcal{T}\right)$是一个拓扑空间，设$\T_1=\{f(S):S\in\T\}$,则$\T_1$是$\mathrm{Y}$上的一个拓扑
\end{theorem}
\begin{proof}
	（1）
	
	由 $\mathrm{X} \in \T,  \Y=f(\mathrm{X} )$ 得 $\Y\in \T_1$
	
	由 $\varnothing \in \T,  \varnothing=f(\varnothing)$ 得 $\varnothing \in \T_1$
	
	（2）$A_\alpha \in \T$ ，则 $f\left(A_\alpha\right) \in \T_1$ ，那么考虑
	
	$$
	\bigcup_\alpha f\left(A_\alpha\right)=f\left(\bigcup_\alpha A_\alpha\right) \in \mathcal{T}_1
	$$
	
	(3)$A_n \in \T$则$f\left(A_n\right) \in \T_1$考虑
	
	$$\bigcap_{k=1}^n f\left(A_k\right)=f\left(\bigcap_{k=1}^n A_k\right) \in \T_1$$
\end{proof}
